\vsssub
\subsection{~程序设计} \label{run:design}
\vssub

\ws{} 的核心是波浪模式子程序，它可由独立程序壳调用，也可由任何需要动态更新波浪数据的其他程序调用。\ws{} 发布包提供了两个这样的程序（例如 {\code ww3\_shel} 和 {\code ww3\_multi}）。辅助程序包括网格预处理器 {\code ww3\_grid}、用于生成人工初始条件的程序 {\code ww3\_strt}、用于单个 {\code ww3\_shel} 或多网格 {\code ww3\_multi} 应用的通用程序壳、两个输入预处理器（{\code ww3\_prep} 和 {\code ww3\_prnc}），以及用于网格输出数据（{\code ww3\_outf} 和 {\code ww3\_ounf}）和点输出数据（{\code ww3\_outp} 和 {\code ww3\_ounp}）的后处理器。

在本节中，文件名用 {\file file} 字体表示，文件内容用 {\code code} 字体表示，{\sc fortran} 程序元素用 {\F fortran} 字体表示。主波浪模式子程序为 {\F w3wave}。除多网格波浪模式 {\file ww3\_multi} 外，数据文件均以文件扩展名 {\file .ww3} 标识；在 {\file ww3\_multi} 中，文件扩展名标识所选多网格镶嵌中的各个网格。为简单起见，本章统一使用文件扩展名 {\file .ww3}。

包含基本数据流的关系图见图~\ref{fig:run_elements}。该图展示了一个典型工作流，如下所述。网格预处理器 {\code ww3\_grid} 写出模式定义文件 {\file mod\_def.ww3}，其中包含海底和障碍物信息，以及定义物理和数值方法的参数值。波浪模式可采用冷启动或热启动。热启动需要重启动文件 {\file restart.ww3}，该文件可由波浪模式在前一次运行中创建，也可由初始条件程序 {\code ww3\_strt} 创建。若重启动文件不可用，波浪模式将自动初始化。若打开线性增长或谱播种，模式可从平静海面（$H_s=0$）开始；否则初始条件将由基于初始风场的参数化有限风区谱组成（见初始条件程序中的相应选项）。

波浪模式例程（{\F w3wave}）可选地生成最多 9 个重启动文件 {\file restart{\em{n}}.ww3}，其中 {\file{\em{n}}} 表示一位整数。对于望远镜式嵌套应用，波浪模式还可选地从文件 {\file nest.ww3} 读取边界条件，并在 {\file nest{\em{n}}.ww3} 中为后续嵌套运行生成边界条件。此外，模式会将原始数据转储到输出文件 {\file out\_grd.ww3}、{\file out\_pnt.ww3}、{\file track\_o.ww3} 和 {\file partition.ww3}（分别为网格平均波浪参数、指定位置处的谱、沿轨迹的谱和分区波浪数据）。需输出谱的轨迹定义在文件 {\file track\_i.ww3} 中。注意，波浪模式不写标准输出，因为当 \ws{} 作为集成模式的一部分时这样会不方便。相反，它维护自己的日志文件 {\file log.ww3}，并可选地为共享内存版本输出测试文件 {\file test.ww3}，或为分布式内存版本输出 {\file test{\em{nnn}}.ww3}，其中 {\em nnn} 是从 1 开始的处理器编号。最后，还提供多种输出后处理器（原始网格场、点输出和轨迹输出文件的二进制后处理；波浪数据的 NetCDF 和 GRIB(2) 打包；用于后续 GrADS 图形处理的网格和谱数据后处理）。下文给出了所有程序元素及其输入文件的更详细说明。注意，每个例程的源代码都有完整文档。这些文档是关于 \ws{} 的额外信息来源。

\input{zh/run/fig_flow_zh}

\ws{} 专用文件在程序子例程中按名称打开。不过，单元号由用户\footnote{~{\file ww3\_multi} 除外。} 在每个 {\file \.inp} 中定义，从而保证在集成模式中实现时具有最大的灵活性。

除波浪模式子程序外，还提供初始化例程和用于数据同化的接口例程。该例程包含一个通用接口，可提供执行完整谱数据同化所需的所有模式组件。该例程集成在通用波浪模式壳中，该壳可为带或不带数据同化的波浪模式执行时间步管理。程序壳还提供一种简单而灵活的方式，将多种数据提供给数据同化方案。数据同化尚未包含在多网格波浪模式壳中。
